\documentclass[14pt]{article}

\usepackage[utf8]{inputenc}
\usepackage[T1]{fontenc}
\usepackage[english]{babel}
\usepackage[a4paper, total={15cm, 24cm}]{geometry}
\usepackage{amsmath}
\usepackage{amssymb}
\usepackage{amsfonts}
\usepackage{graphicx}
\usepackage[font=small]{caption}
\usepackage{subcaption}
\usepackage[dvipsnames]{xcolor}
\usepackage[colorlinks=true, allcolors=Blue]{hyperref}
\usepackage{float}
\usepackage{booktabs}
\usepackage{longtable}
\usepackage{tabularx}
\usepackage{multirow}
\usepackage{adjustbox}
\usepackage{pdflscape}
\usepackage{makecell}
\usepackage{parskip}
\usepackage{setspace}
\usepackage{xurl}

\setstretch{1.2}

\newcommand{\percent}{\%}
\newcommand{\figplaceholder}[2]{%
  \begin{figure}[H]
    \centering
    \IfFileExists{#1}{%
      \includegraphics[width=0.9\linewidth]{#1}%
    }{%
      \fbox{\parbox[c][4cm][c]{0.85\linewidth}{\centering #2\\\small Expected path: \texttt{\detokenize{#1}}}}%
    }%
    \caption{#2}
  \end{figure}
}

\begin{document}

\begin{titlepage}
  \begin{center}
    \IfFileExists{minilogo.png}{\includegraphics[width=0.6\textwidth]{minilogo.png}\\[2cm]}{}
    \vspace*{\fill}

    {\Huge \textbf{Speech commands classification with Transformers}}\\[0.8cm]
    {\Large Deep learning}\\[0.5cm]

    \vspace*{\fill}

    {\Large Mateusz Chudek}\\
    {\large \href{mailto:mateusz.chudek2.stud@pw.edu.pl}{mateusz.chudek2.stud@pw.edu.pl}}\\[0.5cm]

    {\Large Jakub Stachyra}\\
    {\large \href{mailto:jakub.stachyra.stud@pw.edu.pl}{jakub.stachyra.stud@pw.edu.pl}}\\[1.5cm]

    {\large Version 1.0}\\[0.3cm]
    {\large \today}
  \end{center}
\end{titlepage}

\newpage
\tableofcontents
\newpage

\begin{abstract}
This report describes experiments on speech command classification using recurrent and Transformer-based neural networks. The task is based on the TensorFlow Speech Commands dataset, where ten selected commands must be recognized, while all remaining words are grouped into the \texttt{unknown} class and background noise is represented by \texttt{silence}. The main experiments compare LSTM and Transformer models, tune selected hyperparameters, test a two-stage known/unknown cascade, and evaluate sampling strategies for the strongly imbalanced \texttt{unknown} class. The best results in the final experiments were obtained by a single LSTM model trained on the natural data distribution.
\end{abstract}

\section{Introduction}

The goal of the project is to classify short spoken commands. Each audio recording contains either one of ten target commands, an unrelated word assigned to the \texttt{unknown} class, or a fragment of background noise assigned to \texttt{silence}. The project addresses two main objectives: testing and comparing different neural network architectures, with at least one being a Transformer, and investigating the influence of parameter changes on the obtained results. Additionally, the project explores split known/unknown classification, where the model must distinguish between seen and unseen word categories, as well as physical undersampling and oversampling of the training data to address class imbalance.

The project is implemented in Python, with Jupyter notebooks used for running analyses and experiments. Notebook 01 was used for data exploration and analysis. Notebook 02 was used for baseline and initial sanity checks. Notebooks 03, 04 and 05 contain the main experiments discussed in this report:

\begin{enumerate}
    \item hyperparameter experiments,
    \item split known/unknown classification,
    \item physical undersampling and oversampling of the training data.
\end{enumerate}

\section{Dataset and preprocessing}
Experiments are based on TensorFlow Speech Recognition Challenge. The dataset contains short audio recordings sampled at 16 kHz. The target labels are:
\texttt{yes}, \texttt{no}, \texttt{up}, \texttt{down}, \texttt{left}, \texttt{right}, \texttt{on}, \texttt{off}, \texttt{stop}, \texttt{go}, \texttt{unknown}, and \texttt{silence}.
Only a subset of all original words is treated as target commands. Other spoken words are merged into \texttt{unknown} class. This makes the class both large and dataset disbalanced.
Original data from competition comes in two sets, train and test, but the latter was used for competition purposes and has no labels attached.

\figplaceholder{figures/01_target_label_counts.png}{Total and Unique Counts per Class (Target Labels)}

Audio recordings were converted into Mel spectrogram features. The common feature configuration used 64 Mel bands and an FFT size of 512. Training, validation and test splits were kept separate. In the sampling experiment, only the training split was resampled; validation and test data kept the natural distribution.

\section{Theoretical background}

\subsection{Mel spectrograms}

Neural networks usually do not operate directly on raw waveform samples in this project. Instead, the waveform is transformed into a time-frequency representation. A spectrogram shows how energy is distributed over frequencies over time. The Mel scale groups frequencies according to a perceptual scale closer to human hearing, so the model receives features that are more compact and better aligned with speech perception than raw FFT bins.

\subsection{Long Short-Term Memory networks}

Long Short-Term Memory networks are a recurrent neural network architecture designed to process sequences. Standard recurrent networks often struggle with long-term dependencies because gradients can vanish or explode during training. LSTM cells address this by using gates to control the flow of information.

An LSTM can be thought of as having two kinds of memory: a short-term hidden state and a long-term cell state. At each step, three gates decide what to forget from long-term memory, what new information to write into it, and what part of it to expose as output. This allows important information to be preserved across many steps while irrelevant information is discarded early.

For speech commands, the spectrogram is treated as a sequence of time frames which the LSTM processes one by one, building a representation of the whole utterance. A bidirectional LSTM was used in this project, meaning the sequence is processed both forward and backward. This allows later frames to help interpret earlier ones, which is possible since the full recording is available before classification.

\subsection{Transformers}

The Transformer architecture is based on self-attention. Instead of processing a sequence step by step, the model learns relations between all time frames in parallel. In an encoder-only Transformer, each input frame is embedded, positional information is added, and the resulting sequence is passed through Transformer encoder blocks.

The key operation is scaled dot-product attention:

\begin{equation}
    \text{Attention}(Q,K,V) = \text{softmax}\left(\frac{QK^T}{\sqrt{d_k}}\right)V.
\end{equation}

Multi-head attention repeats this operation in several independent heads, allowing the model to learn different relations between time frames. Each encoder block also contains residual connections, layer normalization and a feed-forward network. For this project the decoder part of the original Transformer is not needed, because the task is classification rather than sequence generation.

\section{Initial baseline findings}

Notebook 02 established that both architectures learn the task and that increasing the training fraction improves generalization. These runs are treated as initial findings rather than a standalone experiment, because their purpose was to validate the pipeline and choose reasonable starting settings.

\begin{table}[H]
\centering
\caption{Initial baseline runs from notebook 02.}
\label{tab:baseline}
\begin{tabular}{lrrrrr}
\toprule
Model & Train fraction & Best epoch & Train acc. [\%] & Val. acc. [\%] & Test acc. [\%] \\
\midrule
LSTM & 0.2 & 10 & 90.96 & 88.76 & 88.67 \\
Transformer & 0.2 & 9 & 92.29 & 89.75 & 89.80 \\
LSTM & 0.5 & 9 & 94.26 & 92.49 & 92.51 \\
Transformer & 0.5 & 9 & 94.47 & 92.55 & 92.58 \\
LSTM & 1 & 10 & 96.36 & 94.56 & 94.65 \\
Transformer & 1 & 10 & 95.81 & 94.21 & 93.90 \\
\bottomrule
\end{tabular}
\end{table}

The baseline confirmed that the pipeline was working correctly. The LSTM was slightly better on the full training split, reaching 94.65\% test accuracy compared with 93.90\% for the Transformer. At smaller training fractions, the difference was smaller.

\begin{figure}[H]
\centering
\begin{subfigure}[b]{0.48\linewidth}
    \centering
    \IfFileExists{../reports/models/small_functional_baseline_lstm_transformer/lstm_train1_val1_test1_lr0_001_seed42/figures/learning_curves.png}{\includegraphics[width=\linewidth]{../reports/models/small_functional_baseline_lstm_transformer/lstm_train1_val1_test1_lr0_001_seed42/figures/learning_curves.png}}{\fbox{\parbox[c][3cm][c]{0.9\linewidth}{\centering LSTM baseline learning curves}}}
    \caption{LSTM}
\end{subfigure}
\hfill
\begin{subfigure}[b]{0.48\linewidth}
    \centering
    \IfFileExists{../reports/models/small_functional_baseline_lstm_transformer/trfm_train1_val1_test1_lr0_001_seed42/figures/learning_curves.png}{\includegraphics[width=\linewidth]{../reports/models/small_functional_baseline_lstm_transformer/trfm_train1_val1_test1_lr0_001_seed42/figures/learning_curves.png}}{\fbox{\parbox[c][3cm][c]{0.9\linewidth}{\centering Transformer baseline learning curves}}}
    \caption{Transformer}
\end{subfigure}
\caption{Learning curves for the full-data baseline runs.}
\label{fig:baseline-learning}
\end{figure}

\section{Conducted experiments}

The main experimental work was divided into three notebooks.

\begin{enumerate}
    \item Notebook 03 tested selected hyperparameters while changing one factor at a time.
    \item Notebook 04 tested whether a separate known/unknown gate improves final classification.
    \item Notebook 05 tested whether physical resampling of the training set improves performance on the natural validation and test distributions.
\end{enumerate}

\section{Results}

\section{Experiment 1: hyperparameter search}

Notebook 03 used a fixed experimental configuration so that the influence of individual hyperparameters could be compared directly. Each run used 20\% of the training data, full validation and test splits, natural sampling, and 10 training epochs. Unless a given study changed it, the baseline configuration used dropout 0.2, learning rate 0.001, weight decay 0.0001, batch size 128, and model capacity 64. The experiments changed one parameter at a time for both LSTM and Transformer models, while keeping the remaining settings unchanged.

\subsection{Results}
\small
\begin{longtable}{lllr}
\caption{Full results of notebook 03 hyperparameter experiments.}
\label{tab:hyperparameter-full}\\
\toprule
Study & Architecture & Value & Test acc. [\%] \\
\midrule
\endfirsthead
\toprule
Study & Architecture & Value & Test acc. [\%] \\
\midrule
\endhead
architecture & lstm & lstm & 87.57 \\
\textbf{architecture} & \textbf{transformer} & \textbf{transformer} & \textbf{87.64} \\
\cmidrule(lr){1-4}
dropout & lstm & 0.0 & 87.83 \\
\textbf{dropout} & \textbf{transformer} & \textbf{0.0} & \textbf{89.27} \\
dropout & lstm & 0.1 & 87.94 \\
dropout & transformer & 0.1 & 87.77 \\
dropout & lstm & 0.2 & 87.57 \\
dropout & transformer & 0.2 & 87.64 \\
dropout & lstm & 0.4 & 86.90 \\
dropout & transformer & 0.4 & 84.76 \\
dropout & lstm & 0.5 & 87.80 \\
dropout & transformer & 0.5 & 81.99 \\
\cmidrule(lr){1-4}
learning\_rate & lstm & 0.0001 & 63.59 \\
learning\_rate & transformer & 0.0001 & 66.71 \\
learning\_rate & lstm & 0.0005 & 85.03 \\
learning\_rate & transformer & 0.0005 & 83.40 \\
learning\_rate & lstm & 0.001 & 87.57 \\
learning\_rate & transformer & 0.001 & 87.64 \\
\textbf{learning\_rate} & \textbf{lstm} & \textbf{0.005} & \textbf{88.92} \\
learning\_rate & transformer & 0.005 & 86.40 \\
\cmidrule(lr){1-4}
\textbf{weight\_decay} & \textbf{lstm} & \textbf{0.0} & \textbf{87.77} \\
weight\_decay & transformer & 0.0 & 87.64 \\
\textbf{weight\_decay} & \textbf{lstm} & \textbf{1e-05} & \textbf{87.77} \\
weight\_decay & transformer & 1e-05 & 87.64 \\
weight\_decay & lstm & 0.0001 & 87.57 \\
weight\_decay & transformer & 0.0001 & 87.64 \\
weight\_decay & lstm & 0.001 & 86.44 \\
weight\_decay & transformer & 0.001 & 87.62 \\
\cmidrule(lr){1-4}
\textbf{batch\_size} & \textbf{lstm} & \textbf{32} & \textbf{89.90} \\
batch\_size & transformer & 32 & 89.08 \\
batch\_size & lstm & 64 & 89.17 \\
batch\_size & transformer & 64 & 87.96 \\
batch\_size & lstm & 128 & 87.57 \\
batch\_size & transformer & 128 & 87.64 \\
batch\_size & lstm & 256 & 85.54 \\
batch\_size & transformer & 256 & 84.24 \\
\cmidrule(lr){1-4}
capacity & lstm & 16 & 71.07 \\
capacity & transformer & 16 & 69.26 \\
capacity & lstm & 32 & 80.20 \\
capacity & transformer & 32 & 78.40 \\
capacity & lstm & 64 & 87.57 \\
capacity & transformer & 64 & 87.64 \\
\textbf{capacity} & \textbf{lstm} & \textbf{128} & \textbf{89.47} \\
capacity & transformer & 128 & 88.69 \\
\bottomrule
\end{longtable}

\subsubsection{Comparison plots}

\figplaceholder{figures/03_hyperparameter_accuracy_comparison.png}{Comparison plot for notebook 03 hyperparameter results}
\figplaceholder{figures/03_best_per_study_comparison.png}{Best validation and test accuracy per hyperparameter study}

\subsubsection{Findings}

The baseline LSTM and Transformer were almost tied in notebook 03. The Transformer reached 87.64\% test accuracy and the LSTM reached 87.57\%. The strongest improvements came from decreasing the batch size and increasing model capacity. Batch size 32 was best for both architectures, while batch size 256 clearly reduced accuracy. Capacity 128 was also best for both architectures.

Learning rate was important. A very small value of 0.0001 led to underfitting for both models. The LSTM reached its best learning-rate result at 0.005, while the Transformer was best at 0.001. Dropout affected the Transformer more strongly than the LSTM. Transformer accuracy dropped for high dropout values, while LSTM was relatively stable.

\subsection{Experiment 2: split known and unknown}

Notebook 04 tested a two-stage cascade. The first model is a binary classifier deciding whether the sample is a known command or \texttt{unknown}. If the sample is predicted as known, the second model classifies it into one of the target commands. This approach was compared with a single multiclass model trained on the same split.

\subsubsection{Results}

\begin{table}[H]
\centering
\caption{Full results of notebook 04 split known/unknown experiment.}
\label{tab:cascade-results}
\begin{adjustbox}{width=\linewidth}
\begin{tabular}{lrrrrrrrr}
\toprule
Architecture & Binary acc. & Single acc. & Cascade acc. & Delta & Single known & Cascade known & Single unknown recall & Cascade unknown recall \\
\midrule
LSTM & 94.13 & 92.33 & 90.71 & -1.62 & 85.08 & 84.42 & 96.70 & 94.49 \\
Transformer & 92.85 & 92.36 & 89.90 & -2.46 & 85.20 & 79.78 & 96.67 & 95.99 \\
\bottomrule
\end{tabular}
\end{adjustbox}
\end{table}

\subsubsection{Comparison plots}

\begin{figure}[H]
\centering
\begin{subfigure}[b]{0.48\linewidth}
    \centering
    \IfFileExists{../reports/04_split_unknown_known_experiment/04_split_unknown_known_lstm/lstm_cascade/figures/single_multiclass_learning_curves.png}{\includegraphics[width=\linewidth]{../reports/04_split_unknown_known_experiment/04_split_unknown_known_lstm/lstm_cascade/figures/single_multiclass_learning_curves.png}}{\fbox{\parbox[c][3cm][c]{0.9\linewidth}{\centering LSTM single-model curves}}}
    \caption{LSTM single model}
\end{subfigure}
\hfill
\begin{subfigure}[b]{0.48\linewidth}
    \centering
    \IfFileExists{../reports/04_split_unknown_known_experiment/04_split_unknown_known_lstm/lstm_cascade/figures/binary_learning_curves.png}{\includegraphics[width=\linewidth]{../reports/04_split_unknown_known_experiment/04_split_unknown_known_lstm/lstm_cascade/figures/binary_learning_curves.png}}{\fbox{\parbox[c][3cm][c]{0.9\linewidth}{\centering LSTM binary-gate curves}}}
    \caption{LSTM binary gate}
\end{subfigure}
\caption{Learning curves for the LSTM single model and binary gate in notebook 04.}
\label{fig:cascade-lstm-learning}
\end{figure}

\figplaceholder{figures/04_single_vs_cascade_comparison.png}{Comparison plot of single multiclass and cascade accuracy}

\subsubsection{Findings}

The cascade did not improve accuracy. For LSTM, the cascade was 1.62 percentage points worse than the single multiclass model. For Transformer, the gap was 2.46 percentage points. The binary gate had high standalone accuracy, but errors from the first stage propagated into the second stage. The Transformer cascade especially reduced known-command accuracy, from 85.20\% to 79.78\%.

\subsection{Experiment 3: sampling strategies}

Notebook 05 tested whether the imbalance caused by the \texttt{unknown} class can be improved by physically changing the training set. Three strategies were compared:

\begin{itemize}
    \item \texttt{natural}: no resampling,
    \item \texttt{unknown\_undersampled}: reduce \texttt{unknown} to the size of the largest non-\texttt{unknown} class,
    \item \texttt{undersampled\_then\_oversampled}: reduce the largest class and oversample smaller classes to the same target size.
\end{itemize}

Validation and test splits were not resampled.

\subsubsection{Results}

\begin{table}[H]
\centering
\caption{Full results of notebook 05 sampling experiment.}
\label{tab:sampling-results}
\begin{adjustbox}{width=\linewidth}
\begin{tabular}{llrrrrr}
\toprule
Strategy & Architecture & Train examples & Unknown train examples & Best epoch & Val. acc. [\%] & Test acc. [\%] \\
\midrule
natural & LSTM & 52667 & 32550 & 5 & 92.32 & 92.66 \\
natural & Transformer & 52667 & 32550 & 4 & 91.47 & 91.56 \\
unknown\_undersampled & LSTM & 22002 & 1885 & 5 & 75.23 & 76.16 \\
unknown\_undersampled & Transformer & 22002 & 1885 & 5 & 76.95 & 77.87 \\
undersampled\_then\_oversampled & LSTM & 45240 & 3770 & 5 & 85.27 & 85.17 \\
undersampled\_then\_oversampled & Transformer & 45240 & 3770 & 5 & 82.42 & 83.15 \\
\bottomrule
\end{tabular}
\end{adjustbox}
\end{table}

\subsubsection{Comparison plots}

\begin{figure}[H]
\centering
\begin{subfigure}[b]{0.48\linewidth}
    \centering
    \IfFileExists{../reports/05_sampling_experiment/05_sampling_natural_lstm/lstm_train1_val1_test1_lr0_001_seed42/figures/learning_curves.png}{\includegraphics[width=\linewidth]{../reports/05_sampling_experiment/05_sampling_natural_lstm/lstm_train1_val1_test1_lr0_001_seed42/figures/learning_curves.png}}{\fbox{\parbox[c][3cm][c]{0.9\linewidth}{\centering LSTM natural sampling curves}}}
    \caption{LSTM natural}
\end{subfigure}
\hfill
\begin{subfigure}[b]{0.48\linewidth}
    \centering
    \IfFileExists{../reports/05_sampling_experiment/05_sampling_natural_transformer/trfm_train1_val1_test1_lr0_001_seed42/figures/learning_curves.png}{\includegraphics[width=\linewidth]{../reports/05_sampling_experiment/05_sampling_natural_transformer/trfm_train1_val1_test1_lr0_001_seed42/figures/learning_curves.png}}{\fbox{\parbox[c][3cm][c]{0.9\linewidth}{\centering Transformer natural sampling curves}}}
    \caption{Transformer natural}
\end{subfigure}
\caption{Learning curves for natural sampling in notebook 05.}
\label{fig:sampling-natural-learning}
\end{figure}

\figplaceholder{figures/05_sampling_strategy_comparison.png}{Comparison plot of sampling strategies}

\subsubsection{Findings}

Natural sampling produced the best result for both architectures. LSTM reached 92.66\% test accuracy and Transformer reached 91.56\%. Undersampling \texttt{unknown} caused the largest degradation, reducing LSTM accuracy to 76.16\% and Transformer accuracy to 77.87\%. The mixed undersampling and oversampling strategy was better than pure undersampling, but still clearly worse than the natural distribution.

This suggests that the \texttt{unknown} class should not be treated as a simple redundant majority class. It contains many different words and acoustic patterns. Removing most of it removes useful information about what non-command speech sounds like.

\section{Discussion}

Across the experiments, the simplest successful setup was usually the strongest one: one multiclass model trained on the natural data distribution. The two-stage cascade added complexity but reduced total accuracy. Physical resampling also reduced accuracy because it damaged the representation of the broad \texttt{unknown} class.

The comparison between LSTM and Transformer was not one-sided. The Transformer was competitive in smaller baseline and hyperparameter runs, but the LSTM achieved the strongest full-data baseline and the best natural-sampling result in notebook 05. For this relatively small speech-command task, recurrence is still a strong inductive bias. The LSTM processes spectrograms as ordered temporal sequences and can model short temporal dependencies efficiently.

The most promising tuning directions from notebook 03 are larger capacity and smaller batch size. However, notebook 03 used only 20\% of the training data, so these conclusions should be validated again on the full training split.

\section{Conclusions}

\begin{enumerate}
    \item The baseline experiments confirmed that both LSTM and Transformer models solve the task well. On the full training split, LSTM reached 94.65\% test accuracy and Transformer reached 93.90\%.
    \item Hyperparameter tuning showed that batch size and model capacity had the clearest positive effect. Batch size 32 and capacity 128 were the best tested values in notebook 03.
    \item The two-stage known/unknown cascade did not improve performance. A single multiclass classifier was better for both architectures.
    \item Natural sampling was better than undersampling or oversampling. The large \texttt{unknown} class contains useful diversity and should not be aggressively reduced.
    \item LSTM should remain a strong baseline for this task. Transformers are competitive, but in these experiments they did not consistently outperform LSTM.
\end{enumerate}

\end{document}